\documentclass[10pt,a4paper]{amsart}
\usepackage[margin=19mm]{geometry}
\usepackage{amsmath,amssymb,mathtools}
\usepackage[scr=rsfs,cal=euler]{mathalfa}
\usepackage{xfrac}
\usepackage[unicode,hidelinks,bookmarks=false]{hyperref}
\newcommand{\ie}{i.e.\ }
\newcommand{\cf}{cf.\ }
\newcommand{\resp}{respectively\ }
\newcommand{\Subsection}{\subsection}
\providecommand{\nobreakdash}{\nobreak\hbox{-}}
\setlength{\emergencystretch}{2em}
\multlinegap0pt
\usepackage{amssymb}
\usepackage{mathtools}
\usepackage{xcolor}
\newtheorem{theorem}{Theorem}
\newtheorem{lemma}{Lemma}
\newtheorem{proposition}{Proposition}
\newcommand{\R}{\mathbb{R}}
\newcommand{\RP}{\mathbb{RP}}
\renewcommand{\S}{\mathbb{S}}
\newcommand{\da}{\,dA_{\widehat{g}}}
\newcommand{\dv}{\,dV_{\widehat{g}}}
\newcommand{\hd}{\widehat{\Delta}}
\newcommand{\hg}{\widehat{g}}
\newcommand{\hm}{\widehat{M}}
\newcommand{\hn}{\widehat{\nabla}}
\newcommand{\hw}{\widehat{W}}
\newcommand{\ls}{\{\rho=r\}}
\newcommand{\sls}{\{\rho\le r\}}
\newcommand{\wh}[1]{\widehat{#1}}
\title{Rigidity and volume pinching for the sharp gradient estimate in positive Ricci curvature}
\author{Junyoung Kim, Jiewon Park}
\date{Source: arXiv:2609.03739v1 (CC BY 4.0)}
\begin{document}
\maketitle

\noindent\emph{Source and licence.} Rigidity and volume pinching for the sharp gradient estimate in positive Ricci curvature. Version arXiv:2609.03739v1 (CC BY 4.0), distributed by arXiv under the Creative Commons Attribution 4.0 International licence, http://creativecommons.org/licenses/by/4.0/, as recorded in the arXiv licence metadata for this exact version and shown on the abstract page https://arxiv.org/abs/2609.03739v1. Full licence notice: \texttt{licenses/arxiv-2609.03739-CC-BY-4.0.txt}.\par\smallskip

\noindent\textit{Editorial scope.} Counted unit: the article's theorem thm1 (source0.tex lines 115--121). The article's complete original proof of it is delivered with it (source0.tex lines 417--497, one \texttt{proof} environment of 81 lines, which the article places away from the statement). No mathematics is written, repaired or re-derived: the statement and the proof are the article's own lines, in the article's own notation, and no retained line is substituted or re-flowed.  Retained uncounted as the source-local context the proof needs: lines 110--112; lines 196--198; lines 219--224; lines 275--277; lines 355--362; lines 385--396.  Nothing was omitted from any retained span, so the package carries no omission record.  No retained line contains a citation, so the wrapper carries no bibliography and no bibliography entry.  No retained line contains a figure environment, an image-inclusion command or drawing code, so no image is delivered and the package declares no asset dependency.  Editorial formatting only, in addition: an AMS \texttt{amsart} preamble that restates the article's own declarations and macros used by the retained lines, namely the environments \texttt{theorem}, \texttt{lemma}, \texttt{proposition} on separate counters, together with the standard packages those lines need.  The retained environments print in this document as Theorem 1, Lemma 1, Proposition 1 in order of appearance, whereas the article prints the same items with the numbering of its own full document.  The complete native source of the article as distributed in the arXiv e-print for this exact version is bundled unchanged as \texttt{sources/209313/source0.tex}.
\begin{equation}\label{eq:blow-up-data}
\hg=G_q^2g\quad\text{on $\hm=M\setminus\{q\}$}\quad\text{and}\quad \rho=\left(\frac{G_q}{G_p}\right)^{1/2}.
\end{equation}

\begin{equation}\label{eq:lem2.1}
\hd\rho=3\rho^{-1}|\hn\rho|^2_{\hg},\qquad \hd\rho^2=8|\hn\rho|_{\hg}^2.
\end{equation}

\begin{lemma}\label{lem:Kato}
Define $ V:=G_q^{-1}|\hw|^2_{\hg}$. Then there exists a universal constant $C>0$ such that
\begin{equation}\label{eq:lem2.2}
\hd V\ge-C|\hw|_{\hg}V.
\end{equation} 
\end{lemma}

\begin{equation}\label{eq:assumption sec1}
|\hn\rho|_{\hg}\ge\delta\quad\text{on $M\setminus\{p,q\}$.}
\end{equation}

\begin{proposition}[Asymptotics of $F$ near $0$]\label{prop:asymptotic of F near 0}
As $r\rightarrow0$, we have
\begin{align}
    \label{eq:asymptotic of F near 0}F(r)&=\frac{\omega_3}{4}\gamma^{-4}|W(p)|^2_g\,r^4+O(r^5),\\
    \label{eq:asymptotic of F' near 0}F'(r)&=\omega_3\gamma^{-4}|W(p)|_g^2\,r^3+O(r^4).
\end{align}
In particular, $F(r)=O(r^4)$ and $F'(r)=O(r^3)$ as $r\rightarrow0$.
\end{proposition}

\begin{proposition}[Asymptotics of $F$ as $r\rightarrow\infty$]\label{prop:asymptotic of F at infinity}
The integral
\begin{equation*}
F_\infty:=\int_{\hm}V|\hn\rho|_{\hg}^2\dv
\end{equation*}
is finite, and $F(r)\rightarrow F_\infty$ as $r\rightarrow\infty$. More precisely,
\begin{equation}
\label{eq:asymptotic of F at infinity}
F_\infty-F(r)=\frac{\omega_3}{6}\gamma^{-4}|W(q)|_g^2\,r^{-6}+O(r^{-7})\quad\text{as } r\rightarrow\infty.
\end{equation}
In particular, $F_\infty-F(r)=O(r^{-6})$ as $r\rightarrow\infty$. 
\end{proposition}

\begin{theorem}[Four-dimensional Weyl rigidity]\label{thm1}
There exists a universal constant $c>0$ with the following property. In the setting of \eqref{eq:blow-up-data}, suppose that, for some $\delta>0$, 
\begin{equation*}
|\hn\rho|_{\hg}\ge\delta\quad\text{on $M\setminus\{p,q\}$}\quad \text{and}\quad \sup_{\hm}\rho^2|\hw|_{\hg}\le c\delta^2,
\end{equation*}
then $\hw\equiv 0$ on $\hm$, and $(M,g)$ is isometric to the round sphere $(\S^4,g_{\S^4})$.
\end{theorem}

\begin{proof}[Proof of Theorem \ref{thm1}] We recall that
\begin{equation*}
|\hn\rho|_{\hg}\ge\delta\quad\text{on $M\setminus\{p,q\}$}\quad \text{and}\quad \sup_{\hm}\rho^2|\hw|_{\hg}\le c\delta^2
\end{equation*} by assumption. The coarea formula gives 
\begin{equation*}
F'(r)=\int_{\ls}V|\hn\rho|_{\hg}\da.
\end{equation*}
Using the divergence theorem and the first equation in \eqref{eq:lem2.1}, we have
\begin{equation*}
F'(r)=\int_{\sls}\left(\langle\hn V,\hn\rho\rangle+3\rho^{-1}V|\hn\rho|^2_{\hg}\right)\dv.
\end{equation*}
Differentiating by the coarea formula yields 
\begin{equation*}
F''(r)-\frac{3}{r}F'(r)=\int_{\ls}\partial_\nu V\da=\int_{\sls}\hd V\dv.
\end{equation*}
Next, integration by parts and the second equation in \eqref{eq:lem2.1} gives
\begin{align*}
\int_{\sls}\rho^2\hd V\dv&=r^2\int_{\ls}\partial_\nu V\da-\int_{\sls}\langle\hn\rho^2, \hn V\rangle\dv\\
&=r^2\left(F''-\frac{3}{r}F'\right)-2rF'+8F.
\end{align*}
Hence we have
\begin{equation*}
r^2F''(r)-5rF'(r)+8F(r)=\int_{\sls}\rho^2\hd V\dv.
\end{equation*}
Let 
\begin{equation*}
\omega=\sup_{\hm}\rho^2|\hw|_{\hg}.
\end{equation*}
Lemma \ref{lem:Kato}, together with the assumption \eqref{eq:assumption sec1}, yields that 
\begin{align*}
r^2F''-5rF'+8F&\ge-C\int_{\sls}\rho^2|\hw|_{\hg}V\dv\\&\ge -C\omega\int_{\sls}V\dv\\
&\ge-C\delta^{-2}\omega F(r).
\end{align*}
If we let $\Theta=C\delta^{-2}\omega$, then
\begin{equation}\label{eq:ODE}
r^2F''-5rF'+(8+\Theta)F\ge0.
\end{equation}

We now choose $c>0$ so that $0\le \Theta<1$; since $\Theta=C\delta^{-2}\omega$ and $\omega \le c\delta^2$ by hypothesis, it suffices to take $c=1/(2C)$, where $C$ is the universal constant of Lemma \ref{lem:Kato}. If $\Theta=0$, then $\omega=0$ and $\wh{W}\equiv0$, as desired. Thus, we assume $0<\Theta<1$. Let $x\ge y$ be the roots of the homogeneous indicial equation
\begin{equation*}
\lambda^2-6\lambda+(8+\Theta)=0,
\end{equation*}
that is,
\begin{equation*}
x=3+\sqrt{1-\Theta},\qquad y=3-\sqrt{1-\Theta},
\end{equation*}
and $2<y<3<x<4$. If we define
\begin{align*}
h(r)=\frac{F(r)}{r^x},
\end{align*}
then the inequality \eqref{eq:ODE} implies
\begin{equation}\label{eq:ODE2}
\left(\frac{h'}{r^{y-x-1}}\right)'=\left(\frac{rF'-xF}{r^y}\right)'=r^{-y-1}(r^2F''-5rF'+(8+\Theta)F)\ge0.
\end{equation}
By Proposition \ref{prop:asymptotic of F near 0}, we have
\begin{equation*}
\frac{rF'-xF}{r^y}=O(r^{4-y})\rightarrow0\quad\text{as } r\rightarrow0.
\end{equation*}
It then follows from \eqref{eq:ODE2} that $h'(r)\ge0$ for every $r>0$. On the other hand, Proposition \ref{prop:asymptotic of F at infinity} gives 
\begin{align*}
h(r)=\frac{F(r)}{r^x}\rightarrow0\quad\text{as } r\rightarrow\infty. 
\end{align*}
Since $h$ is non-negative and non-decreasing, it must hold that $h\equiv 0$, and hence $F\equiv 0$. By the lower bound \eqref{eq:assumption sec1} and the definition of $V$, we obtain $\hw\equiv0$ on $\hm$.

The conformal invariance of the Weyl tensor implies $W_g\equiv 0$. Since a locally conformally flat Einstein metric has constant sectional curvature, $(M,g)$ has constant sectional curvature 1. Then, by the Killing--Hopf theorem, $(M,g)$ is isometric to a spherical space form $\S^4/\Gamma$ for some finite group $\Gamma\subset \mathrm{O}(5)$ acting freely on $\S^4$.

In dimension four, the spherical space forms are $\S^4$ or $\mathbb{RP}^4$. To exclude the latter case, we claim that $|\hn\rho|_{\hg}=0$ at some points in the blow-up of $\mathbb{RP}^4$. Denote the covering map by $\pi:\S^4\rightarrow\mathbb{RP}^4$. Fix two distinct points $p,q\in\S^4$ and let $[p]=\pi(p),[q]=\pi(q)\in\RP^4$. Then we have
\begin{equation*}
\pi^\ast G_{[p]}^{\RP^4}=G_{p}^{\S^4}+G_{-p}^{\S^4}.
\end{equation*}
Since we know that $G^{\S^4}_p=(2\sin(d(p,\cdot)/2))^{-2}$ and $d(-p,\cdot)=\pi-d(p,\cdot)$, we have
\begin{equation*}
\pi^\ast G_{[p]}^{\RP^4}=\frac{1}{4\sin^2(d(p,\cdot)/2)}+\frac{1}{4\cos^2(d(p,\cdot)/2)}=\frac{1}{\sin^2d(p,\cdot)}.
\end{equation*}
Thus if we set $u=\langle x,p\rangle, v=\langle x,q\rangle$ on $\S^4\subset\R^5$, then $u^2,v^2$ have the same values at antipodal points, so are well-defined on $\RP^4$. Since we have $G_{[p]}^{\RP^4}=\frac{1}{1-u^2}$, it gives
\begin{equation*}
\rho^2=\frac{G_q}{G_p}=\frac{1-u^2}{1-v^2}.
\end{equation*}

Since $G_{[p]}^{\RP^4}=\frac{1}{1-u^2}$ has minimum on the cut locus $\mathrm{Cut}(p)=\{u=0\}\cong\RP^3$, where $d(p,\cdot)=\pi/2$, $\nabla G_{[p]}^{\RP^4}=0$ there. For the same reason, $\nabla G_{[q]}^{\RP^4}=0$ on $\mathrm{Cut}(q)=\{v=0\}$. Hence, we have $\nabla G_{[p]}^{\RP^4}=\nabla G_{[q]}^{\RP^4}=0$ on $\mathrm{Cut}(p)\cap\mathrm{Cut}(q)\cong\RP^2$, $\nabla\rho=0$ and $\hn\rho=0$ there. Therefore, we deduce $\Gamma$ is trivial and $(M,g)$ is isometric to $\S^4$.  
\end{proof}

\end{document}
